% 星群（综述）
% license CCBYSA3
% type Wiki

本文根据 CC-BY-SA 协议转载翻译自维基百科\href{https://en.wikipedia.org/wiki/Asterism_(astronomy)}{相关文章}。

\begin{figure}[ht]
\centering
\includegraphics[width=8cm]{./figures/c47b677da81e5ec4.png}
\caption{一幅恒星图像，其中有一组显眼的蓝白色亮星。这些亮星共同组成一个名为“衣架”（Coathanger）的星群，形状类似衣架，位于狐狸座（Vulpecula）中。} \label{fig_Asteri_1}
\end{figure}
星群（Asterism）是指在天空中可见的恒星排列或星群图案。星群可以是任何被识别出的恒星形状或组合，因此它的概念比正式定义的 88 个星座更为广泛。星座是以星群为基础的，但与星群不同，星座是具有官方边界的明确区域，整体上覆盖了整个天空。

星群的形态从仅由几颗恒星构成的简单图案，到由大量恒星组成、横跨天空大片区域的复杂结构不等。构成星群的恒星可能是肉眼可见的明亮星，也可能是较暗、甚至需要望远镜才能观测到的星体，但它们通常在亮度上彼此相近。那些较大且较亮的星群对熟悉夜空的人而言具有重要的指引作用。

星群中恒星的排列并不一定反映它们之间的物理联系，而往往是由于观测视角造成的视觉效果。例如，“夏季大三角”仅仅是从地球观察到的一个表观组合，成员恒星之间在空间上并无关联；而“猎户座腰带”中的恒星则都属于猎户座 OB1 星协，“北斗七星”中的七颗星有五颗属于大熊座移动星群（Ursa Major Moving Group）。某些具有真实物理关联的星团，如“毕星团”（Hyades）或“昴星团”（Pleiades），既可以被视为独立的星群，也可能同时是更大星群的一部分。
\subsection{星群与星座的背景}
\begin{figure}[ht]
\centering
\includegraphics[width=8cm]{./figures/590b7d57cdce8a34.png}
\caption{NGC 1664 的图像，其星群呈风筝及尾状，被称为“坎布尔之风筝”。} \label{fig_Asteri_2}
\end{figure}
在许多早期文明中，人们常常将恒星以“连点成线”的方式连接起来，形成类似火柴人形状的图案。其中最早的记录之一来自古印度的《吠檀伽·乔提沙》（Vedanga Jyotisha）以及古巴比伦的天文学文献。\(^\text{[citation needed]}\) 不同的文化辨识出不同的星座，但某些显而易见的恒星排列形式，如猎户座（Orion）和天蝎座（Scorpius），往往会在多个文化的星座体系中同时出现。由于任何人都可以随意组合和命名一组恒星，因此在早期并不存在星座（constellation）与星群（asterism）的明确区别。例如，老普林尼（Pliny the Elder）在其著作《自然史》（Naturalis Historia）中提到了七十二个星群。\(^\text{[3]}\)

包含四十八个星座的一般列表很可能始于天文学家喜帕恰斯（Hipparchus，约公元前190年至公元前120年）。由于当时星座仅被认为是由构成该图形的恒星组成，因此总可以利用剩余的恒星，在既有星座之间创造并插入新的星群。\(^\text{[citation needed]}\)

欧洲人向世界其他地区的探索，使他们接触到了此前从未观测到的恒星。两位以显著扩充南天星座数量而闻名的天文学家是约翰·拜耳（Johann Bayer，1572–1625）与尼古拉·路易·德·拉卡伊（Nicolas Louis de Lacaille，1713–1762）。拜耳列出了十二个图形，这些图形由托勒密（Ptolemy）所无法观测到的南天恒星构成；拉卡伊则在接近南天极的区域创建了十四个新的星群。这些提议中的许多后来被正式接受为星座，而其余部分则保留为星群。\(^\text{[citation needed]}\)

1928年，国际天文学联合会（International Astronomical Union, IAU）根据几何边界将天空精确划分为88个官方星座，将其中包含的所有恒星纳入其范围之内。此后任何新选定的恒星组合或旧有的星座体系，通常都被视为星群。然而，关于“星座”（constellation）与“星群”（asterism）这两个术语之间的技术性区别，至今仍然存在一定的模糊性。\(^\text{[citation needed]}\)
\subsection{由一等星组成的星群}
一些星群完全由明亮的一等星构成，勾勒出简单的几何形状。
\begin{itemize}
\item 由天津四（Deneb）、牛郎星（Altair）和织女星（Vega）——即天鹅座$\alpha$、天鹰座$\alpha$与天琴座$\alpha$——组成的“夏季大三角”，在北半球夏季夜空中格外显眼，因为这三颗恒星皆为一等星。夏季大三角的恒星位于银河带上，该带对应着银河系赤道的方向，并指向银心。
\item “冬季大三角”在北半球冬季夜空中可见，由一等星参宿四（Betelgeuse）、天狼星（Sirius）和南河三（Procyon）组成，其中后两颗分别是肉眼可见的第二近与第四近的恒星系统。
\item 更大的“冬季六边形”则包含天空中二十二颗一等星中的七颗：北河二（Pollux）、五车二（Capella）、毕宿五（Aldebaran）、参宿七（Rigel）、天狼星（Sirius）与南河三（Procyon），外围还有二等星北河三（Castor），以及略偏中心的参宿四（Betelgeuse）。加入参宿四后，它又被称为“天之G形”。该星群围绕着银道反中心，同时包含双子座与猎户座等星座。六边形中，毕宿五的背景处为毕星团（Hyades），这是距离地球最近的疏散星团，也是五个一等深空天体之一；在其东北方向还可见昴星团（Pleiades，亦位于金牛座）以及英仙座α星团（Alpha Persei Cluster，最亮恒星为昴宿六 Alcyone 与天船三 Mirfak）。
\item “春季大三角”由大角星（Arcturus）、轩辕十四（Regulus）和角宿一（Spica）组成。
\item “巨钻形”（Great Diamond）由大角星（Arcturus）、角宿一（Spica）、狮子座$\beta$星（Denebola）和猎犬座$\alpha$星（Cor Caroli）组成，其中后两颗并非一等星。从大角星到狮子座$\beta$星的东西连线，与猎犬座$\alpha$星在北方构成一个等边三角形，与角宿一在南方构成另一个等边三角形，这两个三角形合在一起便形成了“钻石形”。该星群的恒星正式分布于牧夫座、室女座、狮子座与猎犬座。
\end{itemize}
其他一些星群部分由多颗一等星构成。
\begin{itemize}
\item “南十字座”包含一等星南十字座$\alpha$星（Acrux）与$\beta$星（Mimosa），位于船底座星云（五个一等深空天体之一）以西，并有一等星半人马座$\alpha$星（Alpha Centauri，距离太阳最近的恒星）和$\beta$星（Beta Centauri）指向十字形，从而将其与较暗、形状相似的星群——如“钻石十字”或“伪十字”——区分开来。
\end{itemize}
其余所有一等星在其所属的星群或星座中皆为唯一的一等星。例如，大犬座南方的船底座星云以东有船底座$\alpha$星老人星（Canopus），其附近是大麦哲伦云（两者皆为一等深空天体）；波江座中的阿彻纳星（Achernar）位于老人星以东；南鱼座中的北落师门（Fomalhaut）位于阿彻纳星以东；而心宿二（Antares）则位于天蝎座中，视觉上靠近银心位置。
\subsection{基于星座的星群}
\begin{itemize}
\item 北斗七星（The Big Dipper），又称“犁”（The Plough）或“查理之车”（Charles’s Wain），由大熊座中最亮的七颗恒星组成。这些恒星勾勒出熊的后半身与被夸张延伸的尾巴，或可视为构成熊头与颈部上轮廓的“柄”。相比之下，小熊座因其尾部更长而几乎不似熊形，因此更广为人知的名称是“小北斗”（The Little Dipper）。
\item 天鹅座中的“北十字星”（The Northern Cross），其竖轴由位于天鹅尾部的天津四（Deneb，天鹅座$\alpha$）延伸至位于喙部的辇道增七（Albireo，天鹅座$\beta$）；横轴则由天鹅座$\varepsilon$星与$\delta$星两翼相连。
\item “南十字星座”（The Southern Cross）本名即为一星群，但现已被正式认定为独立的星座“南十字座”（Crux）。其主要恒星包括$\alpha$、$\beta$、$\gamma$、$\delta$星，有时也包含$\varepsilon$星。早期在拜耳（Bayer）1603年的《天体图集》（Uranometria）中，南十字仅被视为由半人马座后腿的恒星组成的星群，从而缩小了半人马座的范围。这几颗恒星很可能正是老普林尼在《自然史》（Naturalis Historia）中提到的星群“凯撒之座”（Thronos Caesaris）。
\item “鱼钩星”（The Fish Hook）是夏威夷人对天蝎座的传统称呼。如果将天蝎座图中从心宿二（Antares，天蝎座$\alpha$）到天蝎座$\beta$星与$\pi$星的连线改为从$\beta$星经$\delta$星至$\pi$星的一条弧线，则图像更显“J”形弯钩。在武仙座（Hercules）的主图中，再于左右两翼加上竖线连接肢体，便可构成“蝴蝶星群”（Butterfly Asterism）的图案。
\item 牧夫座（Boötes）有时被称为“冰淇淋筒”（Ice Cream Cone），也被称为“风筝”（Kite）。
\item 仙后座（Cassiopeia）的恒星排列成“W”形，因而这一形状常被用作其别称。
\item 飞马座的大四边形（The Great Square of Pegasus）由壁宿一（Markab）、室宿二（Scheat）、壁宿三（Algenib）与壁宿增七（Alpheratz）组成，代表有翼马的躯干。这一星群早在公元前1100至700年间的楔形文字星表《MUL.APIN》中便已被认作名为“ASH.IKU”（意为“田野”）的星座。如今，壁宿增七仅被视为仙女座的一部分，而在古代它同时属于飞马座与仙女座。
\item 室女座中的“室女之盂”（The Bowl of Virgo）由室女座$\beta$、$\gamma$、$\delta$、$\varepsilon$与$\eta$星组成，加上角宿一（Spica），共同形成一个“Y”形的图案。
\end{itemize}

